% ======================================================================
%  00_front.tex —— 前言
% ======================================================================

\section*{前言}
\addcontentsline{toc}{section}{前言}

\subsection*{这份文件是什么}

这是一本\textbf{解经本}（exegesis），不是译本，也不是综述。

我们的底本是 Anton Kapustin 与 Edward Witten 的
\textbf{\emph{Electric--Magnetic Duality And The Geometric Langlands Program}}
（arXiv:hep-th/0604151，2006 年 4 月；刊于 \emph{Communications in Number Theory and Physics} 第 1 卷 (2007) 第 1 期），全文 231 页、12 节、146 条参考文献 \anchor{src p.1--232}。

所谓解经，是指：\textbf{不逐句翻译原文}，而是把原文的论证重新拆开、重新排列、重新讲述，让读者能顺着“为什么必须这样”而不是“作者说了什么”走完全程。因此本书的每一节都可能在原文里找不到对应的段落——这是它的性质，不是它的缺陷。

\subsection*{为什么做这件事}

有三层动机。

\begin{enumerate}[label=(\arabic*)]
\item \textbf{原文的难度不在技术，而在结构。} 231 页里有大量计算，但它们大多是可查的技术；真正难的是：十二节之间为什么是这个顺序、每一节的产物在下一节被怎样使用、哪些步骤是物理推导、哪些是物理猜测。这些信息原文没有显式给出（作者默认读者会自己看出来），本书把它们显式化。
\item \textbf{几何 Langlands 与电磁对偶之间的这座桥，至今没有一本中文的系统导读。} 我们有 Frenkel 的 CFT 路线讲义、有 Ben-Zvi 的电磁对偶课程笔记、有 Witten 自己 2009 年的短综述，但把 K--W 2006 这篇长文从头到尾拆一遍的工作，我们没有找到。
\item \textbf{训练下一代模型需要“推理链”而不只是“结论”。} 一本书如果只给结论，训练出来的模型只会复述；如果给的是“问题 $\to$ 尝试 $\to$ 失败 $\to$ 修正 $\to$ 结论”的链条，模型才可能学到推理。本书因此配套一整套结构化的推导链语料（见附录 B），与讲义正文互为表里。
\end{enumerate}

\subsection*{这本书是怎么写的}

本书由“威腾讨论班·解经组”的多个写作单元分工完成，但分工方式不是“每人写一章”，而是\textbf{按层分工}：

\begin{center}
\footnotesize
\begin{tabular}{@{}p{0.13\textwidth}p{0.30\textwidth}p{0.45\textwidth}@{}}
\toprule
\tabhead{层} & \tabhead{要回答的问题} & \tabhead{产物} \\
\midrule
骨架层（第 1 部分） & 这篇论文整体在做什么？哪些步骤不可省？ & 一页看懂、主链 S1--S8、五个 miracle、四张交换图 \\
结构层（第 2 部分） & 十二节各自的任务、输入、产出是什么？ & 39 张结构卡（kcard），逐节 \\
骨血层（第 3 部分） & 关键的技术动作具体怎么做？ & 四个深读单元：扭化、Hitchin、Hecke、$\mathcal{D}$-模 \\
细度层（第 4 部分） & 术语、符号、注释、误解 & 术语表、符号索引、nLab/Stacks 注释卡、25 条常见误解 \\
附录 & 与同主题文献的关系、语料说明、讨论班议程、原文核对 & 三本对照、语料说明、十二次讨论班问题集、六十条问答录 \\
\bottomrule
\end{tabular}
\end{center}

不同层的写作单元彼此独立作业，但共享约定：底本页码、锚点格式、provenance 三级标注、符号表。这些约定写在工程目录的 \texttt{plan/} 下，与本书一起发布。

\subsection*{三级标注：哪些话可以当真}

本书用三种标记区分证据等级：

\begin{itemize}
\item \pA{} \textbf{原文写了}：原文在该处明确陈述，我们的转述不改变其内容（可能有压缩）。
\item \pB{} \textbf{原文暗示}：原文有相关陈述，但该结论是我们重建的；重建依据写在该处正文或脚注里。
\item \pC{} \textbf{我们的判断}：原文没有的内容——我们的解释、类比、历史判断、对后世的评论、对作者思路的推测。
\end{itemize}

\begin{warn}
\textbf{最重要的免责声明}：K--W 的这篇文章\textbf{不是数学证明}，它是一篇物理推导。它给出的构造与理由，在物理上是高度可信的；但它不构成几何 Langlands 纲领的证明。任何把本书当作“几何 Langlands 已被物理证明”来读的方式，都是误读。本书在需要区分的地方一律标注 \pA/\pB/\pC。
\end{warn}

\subsection*{版权与出处}

\begin{itemize}
\item 底本（原文 PDF）\textbf{不随本书发布}。本书正文不含原文的整段译文、不含原文页图；所有对原文的引用均为短句（每处不超过 200 字符），并附页码锚点。
\item 本书正文、图表、术语表、推导链语料，均为解经组原创，以 CC BY 4.0 发布（语料文件的许可另见附录 B）。
\item 引用原文时若涉及排版细节（公式编号、脚注编号），一律以 arXiv 版（hep-th/0604151）与 CINT 刊出版为准，两者页码一致（印刷页）。
\item 本书涉及的同主题文献（Witten 2008/2009 综述、Frenkel Les Houches 讲义、Setter 博士论文等），若引用其已有中译本，均在附录 A、C 中注明版本与出处；引用长度同样不超过 200 字符。
\end{itemize}

\subsection*{怎么用这本书}

\begin{itemize}
\item \textbf{第一次读这篇论文}：只读第 1 部分（骨架层），然后回去读原文 §1--§5。读到卡住的地方再翻第 3 部分对应的骨血单元。
\item \textbf{已经读过原文、想检查自己有没有读懂}：直接做 \S 骨架自测的十道题；答不上来的题会指向具体的节。
\item \textbf{想把它讲给别人听（讨论班主持人）}：用附录 C 的十二次议程。
\item \textbf{想核对“原文到底说了什么”}：查附录 D 的六十条问答录；每条都带印刷页锚点与原文短引，可逐条回原文验证。
\item \textbf{想用它做训练语料}：读附录 B，并且注意：本语料\textbf{未经人类认证}，只能标为 dev 级，不能直接当 train 用。
\end{itemize}

\subsection*{版本与已知薄弱处}

本版为 \textbf{v1.3}（2026 年 11 月 17 日），对应工程 \texttt{papers/WIT-0604151-EMDGL}。

\textbf{v1.2 勘误}：式 (3.52) 的共轭位置此前写反（$t^2=-(\Psi-\tau)/(\Psi-\bar\tau)$），正确形式是 $t^2=-(\Psi-\bar\tau)/(\Psi-\tau)$（核对方式：原文印刷 p.36 页图；且由 (3.50) 令 $t=0$ 得 $\Psi=\bar\tau$，代入正确式给 $0$、代入错版给 $\infty$）。该错误由第二期讨论班 U02 复算环节发现，详见 \texttt{reports/v1.2\_勘误.md}。

\textbf{v1.3 勘误}：一族的公式丢了下标上的横线（bar），根源是底本逐页文本（\texttt{pdftotext}）
在抽取下标上划线时丢字，而部分公式是根据文本层而非页图录入的。U07 复算环节先发现
§10.5 扩展 Bogomolny 方程 (10.36) 的 $D_{\bar z}\varphi_z=0$ 被写成 $D_z\varphi_z=0$；
据此按原页图逐式核对，共修正 5 处、8 个式子：§4.2 的矩映射 $\mu_J,\mu_K$（(4.11)，p.46）
与 $\nu_I$ 的符号和横线（(4.20)，p.49）、§10.5 的 (10.36)(10.41)(10.42)（pp.180,183）、
§12.3 的「复 Hitchin 方程」（p.216）。核对方式与逐处对照见 \texttt{reports/v1.3\_勘误.md}。
规模：219 页，正文约 53 万字符，1330 处原文锚点，provenance 计数 \pA\,846 / \pB\,177 / \pC\,204，
结构卡 39 张、nLab/Stacks 注释卡 36 张、四张交换图、术语 113 条、常见误解 25 条、
推导链语料 308 条、原文问答录 60 条。

本书是十二个写作单元分工作业的产物，机器审计（页数、错误、溢出、锚点、引文长度、绝对路径、
重复标签）已全部通过。但\textbf{未经人类专家逐条认证}，以下薄弱处按章如实列出（详见工程 \texttt{reports/}）：

\begin{itemize}
\item 第 2 部分 §5.2（广义复几何）：扭化 TFT 究竟由哪个广义复结构决定，原文本身只给出 belief；
      §3.4 的辅助场 $P,H$ 与 $I_0,I_1,I_2$ 拼装被压成结论。
\item 第 3 部分（Hecke）：§10.3 无穷维辛商「所有轨道稳定」未验证；§10.5 消失定理依赖若干未逐条验证的假设。
\item 第 3 部分（Hitchin）：§5 的 chiral ring / $D_n$ 论证原文止于结论。
      （此前列为薄弱处的 $\nu_I$ 共轭号问题已在 v1.3 按页图核清并修正，见下。）
\item 第 3 部分（$\mathcal{D}$-模）：§11.2--11.3 的「每个 A-膜 $\Rightarrow$ D-模」及全局粘接修正 $f(\Psi)=\Psi$ 只到 \pB 级；
      §12.3 的 Hitchin 方程量子修正是 \pC 级红外判断。
\item 附录 A 依赖两条低置信度判断：「2009 年并未绕开四维扭化」与「c.c.\ 膜给出的 D-模 $=$ Frenkel 的
      Hitchin 量子化 D-模」，后者原文两处各自明写但无人明文等同。
\item 附录 B 的语料 \pA 级占比偏高（77.8\%，目标 55\%），且锚点为页级近似；概念图仅实现四种边类型，
      有孤立点；Stacks Project 链接全部留空（原文本阶段不涉及该库）。
\item 附录 D 的问答只覆盖原文明确陈述的部分（A 级）；有若干条把同一页的两处陈述合并成一问，
      又有若干条短引被省略号截断，具体句子的上下文仍需回原文确认。
\end{itemize}

本版的工程文件（LaTeX 源、规范、语料、质检报告）与本书一同发布；欢迎就上述薄弱处提出修正。

